\subsection{Определения и базовая модель}

Рассмотрим условную модель, используемую Бюро сертификации искусственного интеллекта.

\defi{defi:агент_кандидат}
{Агент-кандидат}{
    Упорядоченная тройка
    \begin{equation*}
        A \is \tup{p, c, h} \in \Nz^3,
    \end{equation*}
    где $p$~--- число правильно решённых контрольных задач, $c$~--- число проверенных ссылок, а $h$~--- число обнаруженных галлюцинаций.
}[
    \item Тройка $A = \tup{3, 4, 0}$ описывает агента, решившего три задачи, проверившего четыре ссылки и не допустившего обнаруженных галлюцинаций.
    \item Тройка $A = \tup{12, 0, 7}$ описывает производительного, но чрезвычайно самоуверенного агента.
][
    В настоящих системах искусственного интеллекта качество нельзя адекватно описать тремя целыми числами.
    Эта модель используется исключительно для демонстрации возможностей шаблона.
]

\defi{defi:рейтинг_агента}
{Рейтинг агента-кандидата $A = \tup{p, c, h}$}{
    Число
    \begin{equation*}
        R(A) \is 2p + c - 3h.
    \end{equation*}
}

\defi{defi:сертифицированный_агент}
{Сертифицированный агент-кандидат}{
    Агент-кандидат $A = \tup{p, c, h}$, для которого выполнены условия
    \begin{equation*}
        R(A) \ge 10
        \quad \text{и} \quad
        h = 0.
    \end{equation*}
}[
    \item Агент $A = \tup{3, 4, 0}$ сертифицирован, поскольку $R(A) = 10$ и $h = 0$.
    \item Агент $A = \tup{5, 3, 1}$ не сертифицирован, поскольку наличие галлюцинации запрещает допуск даже при $R(A) = 10$.
][
    Условие $h = 0$ введено потому, что высокий суммарный рейтинг не должен компенсировать вымышленную ссылку на <<теорему о бесконечно умном тостере>>.
]


\subsection{Утверждения, леммы и теоремы}

\stat{stat:изменение_рейтинга}
{Изменение рейтинга при изменении показателей}{
    Увеличение $p$ на единицу увеличивает рейтинг на $2$, увеличение $c$ на единицу увеличивает рейтинг на $1$, а увеличение $h$ на единицу уменьшает рейтинг на $3$.
}[
    Результат непосредственно следует из \blink{определения рейтинга}{defi:рейтинг_агента}.
    Действительно, коэффициенты при $p$, $c$ и $h$ в выражении $2p + c - 3h$ равны соответственно $2$, $1$ и $-3$.
][
    Последний коэффициент отрицателен, поэтому попытка сгенерировать ещё одну уверенную, но вымышленную ссылку ухудшает положение агента.
]

\lemm{lemm:изменение_рейтинга_после_коррекции}
{Изменение рейтинга после одного раунда коррекции}{
    Пусть $A = \tup{p, c, h}$, где $h \ge 1$, а после одного раунда коррекции показатели агента принимают значения
    \begin{equation*}
        A' \is \tup{p + 2, c + 1, h - 1}.
    \end{equation*}
    Тогда
    \begin{equation*}
        R(A') = R(A) + 8.
    \end{equation*}
}[
    По \blink{определению рейтинга}{defi:рейтинг_агента} имеем:
    \begin{equation*}
        \begin{aligned}
            R(A')
                & = 2(p + 2) + (c + 1) - 3(h - 1) = \\
                & = 2p + c - 3h + 8 = \\
                & = R(A) + 8.
        \end{aligned}
    \end{equation*}
]

\theo{theo:рейтинг_после_нескольких_коррекций}
{Рейтинг после нескольких раундов коррекции}{
    Пусть $A_0 = \tup{p_0, c_0, h_0}$, а каждый раунд коррекции заменяет тройку $\tup{p_n, c_n, h_n}$ тройкой
    \begin{equation*}
        \tup{p_n + 2, c_n + 1, h_n - 1}.
    \end{equation*}
    Тогда для каждого $n \in \Nz$, удовлетворяющего условию $n \le h_0$, выполнено равенство
    \begin{equation*}
        R(A_n) = R(A_0) + 8n.
    \end{equation*}
}[
    \emp{База индукции:} при $n = 0$ имеем $R(A_0) = R(A_0) + 8 \mul 0$, поэтому требуемое равенство выполнено.

    \emp{Шаг индукции:} предположим, что для некоторого $n < h_0$ выполнено равенство $R(A_n) = R(A_0) + 8n$.
    По \blink{лемме об изменении рейтинга после коррекции}{lemm:изменение_рейтинга_после_коррекции} получаем
    \begin{equation*}
        R(A_{n + 1}) = R(A_n) + 8 = R(A_0) + 8n + 8 = R(A_0) + 8(n + 1).
    \end{equation*}

    Следовательно, по методу математической индукции равенство $R(A_n) = R(A_0) + 8n$ выполнено для всех $n \in \Nz$, при $n \le h_0$.
][
    Ограничение $n \le h_0$ необходимо, поскольку после $h_0$ раундов число галлюцинаций становится равным нулю, а дальнейшее уменьшение этого показателя в модели не определено.
]

\conc{conc:сертификация_после_коррекций}
{Сертификация после устранения всех галлюцинаций}{
    Пусть агент $A_0 = \tup{p_0, c_0, h_0}$ проходит $h_0$ раундов коррекции.
    Если
    \begin{equation*}
        R(A_0) + 8h_0 \ge 10,
    \end{equation*}
    то полученный агент $A_{h_0}$ является сертифицированным.
}[
    После $h_0$ раундов число галлюцинаций становится равным нулю.
    По \blink{теореме о рейтинге после нескольких коррекций}{theo:рейтинг_после_нескольких_коррекций} рейтинг полученного агента равен $R(A_0) + 8h_0$ и потому не меньше $10$.
    Следовательно, выполнены оба условия из \blink{определения сертифицированного агента}{defi:сертифицированный_агент}.
]


\subsection{Примеры и замечания}

\exam{exam:сертификация_агента_a17}
{Сертификация агента A-17}{
    Агент A-17 имеет начальные показатели
    \begin{equation*}
        A_0 = \tup{1, 2, 2}.
    \end{equation*}
    Определить, будет ли агент сертифицирован после двух раундов коррекции.
}[
    Начальный рейтинг агента равен
    \begin{equation*}
        R(A_0) = 2 \mul 1 + 2 - 3 \mul 2 = -2.
    \end{equation*}
    После двух раундов коррекции получаем
    \begin{equation*}
        A_2 = \tup{5, 4, 0}.
    \end{equation*}
    По \blink{теореме о рейтинге после нескольких коррекций}{theo:рейтинг_после_нескольких_коррекций} имеем
    \begin{equation*}
        R(A_2) = R(A_0) + 8 \mul 2 = -2 + 16 = 14.
    \end{equation*}
    Число галлюцинаций равно нулю, а рейтинг не меньше $10$, поэтому агент A-17 будет сертифицирован.
][
    После сертификации агенту всё ещё не разрешается самостоятельно изменять собственные критерии сертификации.
]

\note{note:ограниченность_условной_модели}
{Ограниченность условной модели}{
    Введённый рейтинг предназначен только для демонстрации математических блоков.
    Настоящая оценка системы искусственного интеллекта требует анализа качества данных, устойчивости, интерпретируемости, безопасности, воспроизводимости, соблюдения ограничений и множества других факторов.
}


\subsection{Ненумерованные блоки}

Звёздная форма создаёт ненумерованный блок и не увеличивает общий счётчик.

\note*{}
{Почему бюро не учитывает скорость ответа}{
    Агент, мгновенно выдающий неправильный ответ, не становится полезнее агента, которому для того же результата потребовалась минута.
}

\exam*{}
{Пример без ссылки}{
    Агент правильно решил задачу, но сослался на несуществующую статью.
    Следует ли считать ответ полностью корректным?
}[
    Нет.
    Правильный результат не делает вымышленный источник существующим.
]

На ненумерованные блоки не ставятся удалённые ссылки, поэтому их аргумент метки оставляется пустым.


\subsection{Дополнительный материал}

Параметр \verb|optional| отмечает дополнительным весь блок.

\theo[optional]{theo:число_двоичных_протоколов}
{Число двоичных диагностических протоколов}{
    Число последовательностей длины $n \in \Nz$, каждый элемент которых равен $0$ или $1$, составляет $2^n$.
}[
    На каждой из $n$ позиций независимо выбирается одно из двух значений.
    По правилу произведения число возможных последовательностей равно
    \begin{equation*}
        \underbrace{2 \mul 2 \mul \ldots \mul 2}_{n\text{ множителей}} = 2^n.
    \end{equation*}
][
    В бюро значение $1$ означает <<проверка пройдена>>, а значение $0$~--- <<проверка поручена следующей версии модели>>.
]

Параметр \verb|optional=proof| оставляет формулировку обязательной, но отмечает доказательство как дополнительное.

\theo[optional=proof]{theo:повторяющаяся_метка}
{Повторяющаяся метка класса}{
    Если $n + 1$ объектов распределены по $n$ классам, то хотя бы два объекта получат одну и ту же метку класса.
}[
    Предположим противное: каждый класс содержит не более одного объекта.
    Тогда во всех $n$ классах находится не более $n$ объектов, что противоречит наличию $n + 1$ объектов.
]

Параметр \verb|optional={proof,note}| одновременно отмечает дополнительными доказательство и присоединённое замечание.

\theo[optional={proof,note}]{theo:сумма_нечетных_чисел}
{Сумма первых нечётных чисел}{
    Для каждого $n \in \N$ выполнено равенство
    \begin{equation*}
        1 + 3 + \ldots + (2n - 1) = n^2.
    \end{equation*}
}[
    Имеем:
    \begin{equation*}
        \sum_{k = 1}^{n} (2k - 1)
        =
        2 \sum_{k = 1}^{n} k - n
        =
        2 \mul \frac{n(n + 1)}{2} - n
        =
        n^2.
    \end{equation*}
][
    Геометрически равенство можно представить как последовательное достраивание квадратов нечётным числом новых клеток.
    Сотрудники бюро называют эти клетки парковочными местами для очень маленьких роботов.
]

\note*[optional]{}
{Дополнительная историческая справка}{
    Первые версии бюро пытались оценивать интеллект по количеству употреблений слова <<очевидно>>.
    Эксперимент был прекращён после того, как система научилась вставлять это слово перед каждым неверным ответом.
}